\documentclass[10pt]{article}
\usepackage[margin=0.5in]{geometry}
\usepackage{booktabs}
\usepackage{enumitem}
\usepackage[T1]{fontenc}
\usepackage{lmodern}
\setlist{nosep,leftmargin=1.2em}
\setlength{\parindent}{0pt}
\setlength{\parskip}{2pt}
\pagestyle{empty}
\newcommand{\sect}[1]{\vspace{1pt}{\large\bfseries #1}\par}

\begin{document}

\begin{center}
{\Large\bfseries Esch384 Hash on FPGA: Design Document}\\[2pt]
ECE 545 Final Project \quad\textbullet\quad Henry Su, Nick Ngo \quad\textbullet\quad November 10, 2026
\end{center}

\sect{Goal}
Build Esch384 in hardware on an FPGA. Esch384 is a hash function from the
SPARKLE family, a finalist in NIST's Lightweight Cryptography competition.
It takes a message of any length and gives a 384-bit (48-byte) digest.
It is a \emph{sponge}: the message goes in 16 bytes at a time (rate 128 bits),
and the rest of the 512-bit state (capacity 384 bits) is never touched directly.

\sect{Design strategy}
We build a simple \emph{iterative} design first: one permutation step per clock cycle.
We make it fully correct before we try to make it smaller.
Hashing a message of $n$ blocks takes $8n+20$ steps (8 per block, 4 extra on the last block,
and 16 to output the 3 digest pieces).

\sect{Components}
{\small\begin{tabular}{@{}ll@{}}
\toprule
Block & What it does \\
\midrule
Alzette & Mixing box on two 32-bit words: add, rotate, XOR (4 rounds). 8 copies, one per branch. \\
Linear layer & Mixes the 8 branches together and swaps their order. \\
SPARKLE512 & One step = constant + 8 Alzettes + linear layer; repeated 8 or 12 times. \\
Datapath & 512-bit state register, message padding and injection, digest output. \\
Controller & State machine: take block, run steps, output 3 digest pieces; valid/ready handshake. \\
Top level & Datapath + controller, tested against all known-answer tests. \\
\bottomrule
\end{tabular}}

\sect{Tools and libraries}
SystemVerilog (kept close to Verilog-2001); Xilinx Vivado 2026.1 for synthesis and xsim for
simulation (Questa as backup); target board Arty A7 (Artix-7, \texttt{xc7a100tcsg324-1}).
Reference: the official NIST C code for Esch384, plus Python scripts that create test vectors and
compare outputs. No outside hardware libraries (IP cores). Code is on GitHub.
\textbf{AI tool:} Claude Code (Anthropic). It drafts RTL, testbenches and scripts, and helps with
debugging; we run every test ourselves before accepting its work. Our proposal named GPT 6 Astra;
we switched to Claude Code. All AI tool use is logged in the repo (\texttt{docs/ai\_log.md}).

\sect{Optimization target: area, and planned optimizations}
Lightweight cryptography is about small area, so area is our target. For every version we report
LUTs, flip-flops, maximum clock speed, latency and throughput.
\begin{itemize}
\item \textbf{Serialize the Alzettes:} use 4, 2 or 1 Alzette box instead of 8, and take more cycles per step,
like the 32-bit serial AES datapath in lecture 13 (Aysu et al., ESL'14).
\item Share adders, reuse registers, try a narrower datapath. One git branch per idea; all tests re-run after every change.
\end{itemize}

\sect{Risk management}
\begin{itemize}
\item \textbf{Wrong results:} every block is tested on its own against the official C code, and we break a copy
on purpose to make sure the test catches it. For bugs, the reference prints the state after every step, so we can find the first step that differs.
\item \textbf{Schedule:} tasks are small and split between us. The working version gets a git tag
before any optimization, so we can always go back to it.
\end{itemize}

\sect{Progress and plan}
{\small\begin{tabular}{@{}llll@{}}
\toprule
Item & Owner & Status / due & Result \\
\midrule
Setup: tools, repo, course notes & Henry & Done & Vivado/xsim 2026.1 working \\
Reference model and test vectors & Henry & Done & Matches all 1025 official test vectors \\
Alzette & Henry & Done & 64/64 tests pass; 254 LUTs \\
Linear layer & Henry & Done & 64/64 tests pass; 320 LUTs \\
SPARKLE512 permutation & Henry & Done & 71/71 tests pass; 2023 LUTs, 521 FFs; max clock 76 MHz \\
Datapath and controller & Nick & Nov 17 & Specs and test plan written \\
Full hash, all test vectors & Nick & Nov 21 & -- \\
Synthesis: area, speed, throughput & Henry & Nov 24 & -- \\
Area optimization & Both & Dec 1 & -- \\
Final presentation and report & Both & Dec 9 & -- \\
\bottomrule
\end{tabular}}

{\small LUT/FF counts are post-synthesis (Vivado, module on its own); 76 MHz is after place and route. 2023 is less than
$8\times254+320$ because inside the permutation the constants are fixed, so Vivado simplifies them
and packs the linear-layer XORs into the same LUTs as the Alzette logic.
SPARKLE512 takes 9 / 13 cycles: 1 load cycle + 8 or 12 steps.}

\sect{Challenges and how we solved them}
\begin{itemize}
\item \textbf{Byte and word order were unclear in the spec.} We checked them against the official C code.
\item \textbf{The simulator (xsim) crashed} on one testbench line. We rewrote the line a simpler way.
\end{itemize}

\end{document}
